\chapter{Numerical Methods for Pricing}
\label{ch:numerical_methods}
\chaptermark{Numerical methods for pricing}

Most derivatives lack closed-form prices.  Black--Scholes
(\cref{ch:black_scholes}) is the famous exception; American options,
path-dependent exotics, multi-asset baskets, callable bonds, CDOs, and
structured notes must be priced numerically.  This chapter covers the
three pillars a Strat uses daily --- \textbf{Monte Carlo} (with
variance reduction), \textbf{finite-difference PDE}, and
\textbf{lattices (trees)} --- plus the \textbf{Longstaff--Schwartz}
algorithm, which converts the American-option optimal-stopping problem
into a sequence of linear regressions.  The aim is to let you
\emph{read} production code, understand its failure modes, and pick
the right method for a given product; a decision table at the end
summarises the choice.

\begin{backgroundbox}[Prerequisites]
From earlier chapters:
\begin{itemize}
  \item \Cref{ch:mean_reversion_ou}: stochastic differential equations
    (SDEs), Brownian motion, the Euler--Maruyama scheme.
  \item \Cref{ch:black_scholes}: the Black--Scholes PDE
    $\partial_t V + \tfrac12 \volat^2 \Spx^2 \partial_{SS} V
     + \rrf \Spx \partial_S V - \rrf V = 0$,
    and the closed-form European call
    $\Call = \Spx \CDFN(d_1) - \Kstrike e^{-\rrf \Tmat} \CDFN(d_2)$.
  \item \Cref{ch:greeks_hedging}: Greeks as partial derivatives of the
    price; why accuracy in Greeks (not just price) matters.
  \item \Cref{ch:fixed_income_foundations}: discount factors
    $D(t) = e^{-\rrf t}$ and the PV-of-future-cashflow mantra.
\end{itemize}
Linear algebra (matrix--vector products, tridiagonal solves) is
assumed.  Numerical PDE and lattice sections use only elementary
Taylor expansions --- if you have seen the heat equation and a forward
difference, you have enough.
\end{backgroundbox}


%% ================================================================
\section{Monte Carlo pricing}
\label{sec:mc_pricing}
%% ================================================================

Risk-neutral pricing (\cref{ch:black_scholes}) gives the European-claim
price
\begin{equation}
\label{eq:rn_pricing}
  V_0 \;=\; e^{-\rrf \Tmat}\, \E^{\Qm}\!\left[ g(\Spx_\Tmat) \right],
\end{equation}
under the risk-neutral measure $\Qm$.  Monte Carlo takes this at
face value: simulate thousands of random price paths under $\Qm$,
evaluate the payoff at each terminal (or along each path), average,
and discount.  When the payoff depends on the whole path --- Asian
averages, barrier knock-outs, lookback maxima --- no closed form
exists, and this path-by-path averaging is often the only tool that
works.

\subsection{The algorithm}

For a European payoff on a single asset under geometric Brownian
motion (GBM),
\begin{equation}
  \Spx_\Tmat \;=\; \Spx_0 \exp\!\left[
    \left(\rrf - \tfrac12 \volat^2\right)\Tmat
    + \volat \sqrt{\Tmat}\, Z \right],
  \quad Z \sim \Normal(0,1),
\end{equation}
the procedure is four lines:
\begin{enumerate}
  \item Draw $Z_1, \ldots, Z_N \iid \Normal(0,1)$.
  \item Set $\Spx_\Tmat^{(k)} = \Spx_0 \exp[(\rrf - \tfrac12\volat^2)\Tmat + \volat\sqrt{\Tmat}\,Z_k]$.
  \item Compute payoff $g_k = g(\Spx_\Tmat^{(k)})$.
  \item Estimate $\hat V_0 = e^{-\rrf\Tmat}\, \tfrac1N \sum_k g_k$.
\end{enumerate}

For path-dependent payoffs (Asian, lookback, barrier), step~2
simulates the whole path by Euler--Maruyama on $0 = t_0 < t_1 < \cdots
< t_M = \Tmat$:
\begin{equation}
  \Spx_{t_{n+1}} = \Spx_{t_n} \exp\!\left[
    (\rrf - \tfrac12 \volat^2)\Delta t + \volat\sqrt{\Delta t}\, Z_{n,k}
  \right].
\end{equation}
The exponential form is exact for GBM and removes time-step bias.

\subsection{Convergence and standard error}

The MC estimator $\hat V_0$ is unbiased, with standard error
\begin{equation}
  \mathrm{SE}(\hat V_0)
  \;=\; e^{-\rrf\Tmat}\, \frac{\sigma_g}{\sqrt{N}},
  \qquad \sigma_g^2 \defeq \Var\!\left[ g(\Spx_\Tmat) \right].
\end{equation}
Halving error quadruples path count --- the defining number about MC:

\begin{keyfactbox}[Monte Carlo convergence rate]
The root-mean-square error of a plain-vanilla Monte Carlo estimator
is $O(1/\sqrt{N})$ regardless of dimension.  This
dimension-independence is what makes MC the only practical method for
high-dimensional problems (basket options, CDOs, multi-factor rates),
but $O(1/\sqrt{N})$ also means brute force alone is slow --- hence
variance reduction.
\end{keyfactbox}

\subsection{Pseudo-random versus quasi-random: a brief aside}

The $O(1/\sqrt N)$ rate assumes the $Z_k$ are independent draws.
\emph{Quasi-Monte Carlo} (QMC) replaces pseudo-random uniforms with
a low-discrepancy sequence (Sobol, Halton, Faure) that covers
$[0,1]^d$ more uniformly than random noise.  Its error is
$O(\log^d(N)/N)$ --- effectively $O(1/N)$ --- and dominates MC for
moderate dimensions ($d \lesssim 50$).  Sobol with Brownian-bridge
construction is the production choice for low-to-moderate dimensional
European exotics.  A naive sequential time-grid fill concentrates the
first few coordinates on ``cost-heavy'' early steps and destroys the
QMC benefit, so paths are always constructed by Brownian bridge or
principal-components decomposition.

\subsection{Worked example: European call by MC}

Let $\Spx_0 = 100$, $\Kstrike = 100$, $\rrf = 5\%$, $\volat = 20\%$,
$\Tmat = 1$ year.  Black--Scholes gives
\begin{equation}
  \Call^{\mathrm{BS}} \;=\; 100\,\CDFN(d_1) - 100\,e^{-0.05}\CDFN(d_2)
  \;=\; 10.4506.
\end{equation}
Drawing $N = 10^6$ standard normals (seed 42) yields
\begin{equation}
  \hat{\Call}^{\mathrm{MC}} = 10.4342,
  \quad \mathrm{SE} = 0.0147,
  \quad |\hat{\Call} - \Call^{\mathrm{BS}}| = 0.0164.
\end{equation}
The absolute error is within $1.1\,\mathrm{SE}$ of the true price,
consistent with the estimator's Gaussian sampling
distribution~\citeGlasserman.  Pushing the error below one cent needs
$(0.0147/0.005)^2 \approx 9$ times more paths --- where variance
reduction earns its keep.

The $O(1/\sqrt N)$ rate is visible in a convergence table for the
same setup (seed 42):

\begin{center}
\small
\renewcommand{\arraystretch}{1.15}
\begin{tabular}{r|rrr}
\toprule
$N$ (paths) & $\hat{\Call}^{\mathrm{MC}}$ & $\mathrm{SE}$ & $|\hat{\Call}-\Call^{\mathrm{BS}}|$ \\
\midrule
$10^3$  & $10.25$ & $0.46$ & $0.20$ \\
$10^4$  & $10.47$ & $0.15$ & $0.02$ \\
$10^5$  & $10.44$ & $0.046$ & $0.009$ \\
$10^6$  & $10.434$ & $0.015$ & $0.016$ \\
\bottomrule
\end{tabular}
\end{center}

Each tenfold path increase shrinks SE by $\sqrt{10} \approx 3.16$ ---
one significant digit.  Six-digit accuracy from $10^6$ paths would
need $10^{12}$ paths (weeks of compute); variance reduction replaces
that brute force.

\begin{intuitionbox}
Monte Carlo \emph{is} the Law of Large Numbers applied to
\eqref{eq:rn_pricing}: any European payoff is priced by replacing the
integral against an unknown density with an empirical average along a
simulation.  Variance reduction, Longstaff--Schwartz, and antithetics
all make that empirical average a better estimator.
\end{intuitionbox}


%% ================================================================
\section{Variance reduction}
\label{sec:variance_reduction}
%% ================================================================

Since MC error scales as $1/\sqrt N$, cutting \emph{variance} by a
factor $k$ is equivalent to running $k$ times more paths for free.
Three techniques dominate production: antithetic variates, control
variates, and importance sampling, each exploiting payoff structure
brute force ignores.

\subsection{Antithetic variates}

If $Z \sim \Normal(0,1)$, so does $-Z$.  Pairing each draw with its
sign-flipped twin gives the estimator
\begin{equation}
  \hat V^{\mathrm{AV}}
  \;=\; e^{-\rrf\Tmat}\,
        \frac{1}{N} \sum_{k=1}^{N/2}
        \tfrac12 \bigl[ g(\Spx_\Tmat(Z_k)) + g(\Spx_\Tmat(-Z_k)) \bigr],
\end{equation}
which uses $N$ payoff evaluations organised into $N/2$ correlated
pairs.  Its variance is
\begin{equation}
  \Var(\hat V^{\mathrm{AV}})
  \;=\; \frac{1}{2N}
        \bigl[ \Var(g(Z)) + \Cov(g(Z), g(-Z)) \bigr].
\end{equation}
If $g$ is monotone in $Z$ (e.g.\ a call on GBM: higher $Z$ $\to$
higher $\Spx_\Tmat$ $\to$ higher payoff), then $\Cov(g(Z), g(-Z)) < 0$
and variance drops.  For antithetics to help, the payoff must be
(anti-)symmetric enough in the driving shock.

\textbf{Empirical check.}  For the European call above, plain MC with
$N = 10^6$ gives $\mathrm{SE} = 0.0147$; the same paths as $5\times
10^5$ antithetic pairs give $\mathrm{SE} = 0.0104$, a $\sim 2\times$
variance cut.  For a linear payoff the reduction is infinite; the
ITM cut-off drags the gain down.

\subsection{Control variates}

Want $\mu_Y = \E[Y]$ for a complicated payoff $Y$, and have a
correlated $X$ with known $\mu_X = \E[X]$ in closed form.  The
control-variate estimator is
\begin{equation}
\label{eq:cv_estimator}
  \hat\mu_Y^{\mathrm{CV}}
  \;=\; \bar Y - b\,(\bar X - \mu_X),
  \quad \text{with variance}\quad
  \Var(\hat\mu_Y^{\mathrm{CV}})
  = \tfrac1N \bigl[ \Var(Y) - 2b\,\Cov(X,Y) + b^2 \Var(X) \bigr].
\end{equation}
Minimising over $b$ gives $b^\star = \Cov(X,Y)/\Var(X)$, at which
\begin{equation}
  \Var(\hat\mu_Y^{\mathrm{CV}})
  \;=\; \tfrac1N \Var(Y) \,(1 - \rho_{XY}^2),
\end{equation}
where $\rho_{XY}$ is the correlation between $X$ and $Y$.  Variance
reduction is $(1-\rho^2)^{-1}$; $\rho = 0.9$ already buys $\sim
5\times$.

\textbf{Worked example: lookback call controlled by vanilla.}
A floating-strike lookback call on GBM pays $(\max_{0\le t\le
\Tmat}\Spx_t - \Kstrike)^+$.  With the same parameters, $N=10^5$ paths,
$M=100$ Euler steps, use the terminal vanilla payoff $(\Spx_\Tmat -
\Kstrike)^+$ as control $X$ (expectation: BS price $10.4506$).
Empirically,
\begin{itemize}
  \item plain MC: $\widehat{LB} = 17.867$, $\mathrm{SE} = 0.0483$;
  \item with CV: $\widehat{LB} = 17.904$, $\mathrm{SE} = 0.0209$;
  \item $\rho(Y,X) = 0.901$, $b^\star = 0.938$, variance ratio
    $\approx 5.3$.
\end{itemize}
The control variate cut the path count $5\times$ for the same
accuracy at negligible extra cost.

\subsection{Variance reduction cheatsheet}

Side-by-side comparison on the same ATM call ($\Spx_0=100, \Kstrike=100,
\rrf=5\%, \volat=20\%, \Tmat=1$, $N = 10^5$):

\begin{center}
\small
\renewcommand{\arraystretch}{1.15}
\begin{tabular}{p{0.26\linewidth} | r | p{0.42\linewidth}}
\toprule
\textbf{Method} & \textbf{SE} & \textbf{Comment} \\
\midrule
Plain MC           & $0.046$
  & Baseline, $N=10^5$ paths. \\
Antithetic         & $0.033$
  & $\approx 2\times$ variance cut; free. \\
Control variate (geom Asian) & $0.002$
  & $\approx 500\times$ variance cut for arithmetic-Asian; massive for products with tight-correlated closed-form cousins. \\
Importance sampling, $\mu=0.6$  & $0.010$
  & Useful only deep OTM; neutral here. \\
Quasi-MC (Sobol)   & $0.006$
  & $\approx 60\times$ variance cut; error scales as $\sim 1/N$ rather than $1/\sqrt N$. \\
\bottomrule
\end{tabular}
\end{center}

The standout is always a good control variate.  Antithetic is a
$2\times$ freebie, importance sampling shines only in the tail, and
QMC is near-universal but needs careful path construction.  Production
libraries apply antithetics and Sobol by default and add a product-
specific control variate when one exists.

\subsection{Stratified sampling and Latin hypercube}

Stratified sampling partitions the sample space into $K$ equal-
probability strata and draws a fixed number of samples from each.
For GBM, split the $Z$ axis into bins $\CDFN^{-1}(k/K)$ and draw
$N/K$ samples per bin.  The variance is
\begin{equation}
  \Var(\hat V^{\mathrm{strat}})
  \;=\; \frac{1}{N} \sum_{k=1}^K \pi_k \Var(g \mid \text{stratum } k)
  \;\leq\; \frac{1}{N} \Var(g),
\end{equation}
with equality only if $g$ is constant within each stratum.  For a
near-ATM call, stratifying on the terminal shock gives $2$--$5\times$
variance reduction.  Latin hypercube extends this to multiple
dimensions by stratifying each coordinate and reshuffling.

\textbf{Stacking.}  Production pricers stack techniques: a typical
autocallable pricer uses Sobol numbers + Brownian-bridge paths +
antithetics + control variate on the digital component.  The
reductions compound multiplicatively and can exceed $100\times$,
turning a three-hour MC into two minutes.

\subsection{Importance sampling}

For deep-OTM options, most paths contribute zero; the signal lives
in the rare tail.  Importance sampling shifts the sampling measure so
tail paths become typical, then re-weights via the Radon--Nikodym
derivative:
\begin{equation}
  \E^{\Qm}[g(\Spx_\Tmat)]
  \;=\; \E^{\widetilde\Qm}\!\left[ g(\Spx_\Tmat)\,
        \frac{d\Qm}{d\widetilde\Qm} \right],
\end{equation}
the trick is to pick $\widetilde\Qm$ so the integrand is as flat as
possible.  For GBM, drift the Brownian by $\mu \Tmat$: sample $Z \sim
\Normal(\mu, 1)$ and multiply the payoff by $\exp(-\mu Z + \tfrac12
\mu^2)$.  Choosing $\mu$ so the typical path lands near the strike
can reduce variance by orders of magnitude for deep-OTM options.  See
\citeGlasserman, chapter~4, for the rate-function heuristic.

\begin{keyfactbox}[Three variance reduction techniques]
\begin{enumerate}
  \item \textbf{Antithetic variates.}  Reuse $Z_k$ and $-Z_k$.  Free
    for any payoff with monotone sensitivity to the driving shock.
  \item \textbf{Control variates.}  Subtract a closed-form-priced
    correlated payoff.  Variance reduction $= (1-\rho^2)^{-1}$; works
    whenever the product has a simpler cousin whose price is known.
  \item \textbf{Importance sampling.}  Sample from a tilted measure
    and re-weight.  Essential for deep OTM, large-loss CVaR, rare-
    default simulation, and any payoff dominated by a thin tail.
\end{enumerate}
\end{keyfactbox}

\begin{intuitionbox}
Variance reduction buys paths.  A $4\times$ cut = $4\times$ fewer
paths for the same accuracy, or the same paths for half the error.
Production pricers stack antithetics, a control variate, and
stratified sampling; the savings compound multiplicatively.
\end{intuitionbox}

\begin{pitfallbox}[Antithetic variates can backfire]
If the payoff is not monotone in the shock --- a straddle, a
butterfly, anything with a kink in both directions --- antithetics can
\emph{increase} variance.  Always measure the post-antithetic standard
error empirically rather than assume the technique helps.
\end{pitfallbox}


%% ================================================================
\section{Longstaff--Schwartz for American options}
\label{sec:longstaff_schwartz}
%% ================================================================

American options (exercisable any time before expiry) are the default
in equity derivatives.  A single-stock American put is worth strictly
more than its European counterpart because of early exercise.  The
pricing problem is \emph{optimal stopping}: the optimal exercise time
depends on the future distribution conditional on the current state
--- exactly what simulations do not give path-by-path.

Before 2001, practical American pricers were limited to lattices and
low-dimensional PDEs.  \citet{longstaffschwartz2001} changed that:
\emph{replace the unknown continuation value with its regression on
the current state}, then compare the regression's prediction to the
immediate-exercise payoff.  The whole method is a sequence of OLS
regressions.

\subsection{The problem statement}

Let $\{\Spx_t\}_{t=0}^{\Tmat}$ be simulated under $\Qm$ with exercise
dates $t_0 < t_1 < \cdots < t_M$.  The American option value is
\begin{equation}
  V_0 \;=\; \sup_{\tau \in \mathcal{T}} \E^{\Qm}\!\left[
            e^{-\rrf \tau} g(\Spx_\tau) \right],
\end{equation}
where the supremum is over stopping times $\tau$ taking values in
$\{t_0, \ldots, t_M\}$.  By dynamic programming,
\begin{equation}
\label{eq:dpeqn}
  V(t_m, \Spx) \;=\; \max\!\left\{
    g(\Spx),\;
    e^{-\rrf \Delta t}\, \E^{\Qm}\!\bigl[ V(t_{m+1}, \Spx_{t_{m+1}})
                                         \mid \Spx_{t_m} = \Spx \bigr]
  \right\}.
\end{equation}
The conditional expectation inside is the \emph{continuation value}
$C(t_m, \Spx)$.  A lattice knows $C$ exactly (it is a two-point
average); a Monte-Carlo simulation does not.

\subsection{The regression idea}

Longstaff--Schwartz approximate $C(t_m, \cdot)$ on the in-the-money
paths by a linear combination of basis functions
$\{\psi_1, \ldots, \psi_L\}$:
\begin{equation}
  C(t_m, \Spx) \;\approx\; \sum_{\ell=1}^{L} \beta_\ell \psi_\ell(\Spx).
\end{equation}
Typical bases are monomials $\{1, \Spx, \Spx^2\}$ or Laguerre/Hermite
polynomials.  The coefficients $\beta_\ell$ come from regressing the
realised discounted future cash flow against $\psi_\ell(\Spx_{t_m})$.

\begin{keyfactbox}[Longstaff--Schwartz algorithm]
\textbf{Backward sweep} for an American put with $M+1$ exercise dates
and $N$ simulated paths:
\begin{enumerate}
  \item At $t_M$: set the realised cash flow on each path to the
    terminal payoff $(\Kstrike - \Spx_{t_M})^+$.
  \item For $m = M-1, M-2, \ldots, 1$:
  \begin{enumerate}
    \item Identify in-the-money paths at $t_m$
      ($\Kstrike > \Spx_{t_m}^{(k)}$).
    \item On those paths, let $Y^{(k)}$ be the realised future cash
      flow discounted back to $t_m$.
    \item Regress $Y$ on $\{\psi_\ell(\Spx_{t_m})\}$ to get
      $\hat C(t_m, \Spx_{t_m}^{(k)})$.
    \item Exercise on path $k$ if $g(\Spx_{t_m}^{(k)}) > \hat C(t_m,
      \Spx_{t_m}^{(k)})$; otherwise carry the existing future cash
      flow.  If exercising, zero out any later CF on that path.
  \end{enumerate}
  \item Discount every path's cash flow back to $t_0$ and average.
\end{enumerate}
The regression uses only in-the-money paths because
out-of-the-money exercise is never optimal: the regression needs to
learn the continuation value exactly where the exercise decision is
non-trivial.
\end{keyfactbox}

\subsection{Worked example: the Longstaff--Schwartz 2001 paths}

We reproduce the original paper's hand-worked example exactly.  Take
$N=8$ paths, $M=3$ annual exercise dates (so four time points
$t=0,1,2,3$), strike $\Kstrike = 1.10$, risk-free rate $\rrf = 6\%$,
discount factor per step $e^{-0.06} = 0.9418$, and basis
$\{1, \Spx, \Spx^2\}$.  The paths are

\begin{center}
\small
\begin{tabular}{c|cccc}
\toprule
path & $t=0$ & $t=1$ & $t=2$ & $t=3$ \\
\midrule
1 & 1.00 & 1.09 & 1.08 & 1.34 \\
2 & 1.00 & 1.16 & 1.26 & 1.54 \\
3 & 1.00 & 1.22 & 1.07 & 1.03 \\
4 & 1.00 & 0.93 & 0.97 & 0.92 \\
5 & 1.00 & 1.11 & 1.56 & 1.52 \\
6 & 1.00 & 0.76 & 0.77 & 0.90 \\
7 & 1.00 & 0.92 & 0.84 & 1.01 \\
8 & 1.00 & 0.88 & 1.22 & 1.34 \\
\bottomrule
\end{tabular}
\end{center}

\textbf{Step 1 (terminal payoff).}
$(\Kstrike - \Spx_3)^+ = (0,\,0,\,0.07,\,0.18,\,0,\,0.20,\,0.09,\,0)$.

\textbf{Step 2 ($t=2$).}  In-the-money paths are $\{1,3,4,6,7\}$ with
$\Spx_2 = (1.08, 1.07, 0.97, 0.77, 0.84)$.  Their discounted future
cash flows are $Y = (0, 0.0659, 0.1695, 0.1884, 0.0848)$ (multiply the
$t=3$ payoff by $e^{-0.06}$).  Regressing $Y$ on $\{1,\Spx,\Spx^2\}$
gives coefficients $\hat\beta = (-1.070,\, 2.983,\, -1.814)$ and
continuation estimates $\hat C = (0.037, 0.046, 0.118, 0.152, 0.156)$.
Immediate exercise pays $\Kstrike - \Spx_2 = (0.02, 0.03, 0.13, 0.33,
0.26)$.  Exercise now on paths $\{4, 6, 7\}$; hold on $\{1, 3\}$.

\textbf{Step 3 ($t=1$).}  ITM paths are $\{1, 4, 6, 7, 8\}$ with
$\Spx_1 = (1.09, 0.93, 0.76, 0.92, 0.88)$.  Discounted future cash
flows (taking into account the new exercise times from $t=2$) are
$Y = (0, 0.1224, 0.3108, 0.2449, 0)$.  Regression gives $\hat\beta =
(2.038,\, -3.335,\, 1.356)$ and $\hat C = (0.013, 0.109, 0.286, 0.117,
0.153)$.  Immediate exercise pays $(0.01, 0.17, 0.34, 0.18, 0.22)$.
Exercise on $\{4, 6, 7, 8\}$; hold on $\{1\}$.

\textbf{Step 4.}  Discount each path's exercise cash flow back to
$t=0$ and average over all 8 paths:
\begin{equation}
  \hat V_0 \;=\; \frac{1}{8}\bigl(
      e^{-0.18}\cdot 0.07
    + e^{-0.06}\cdot 0.17
    + e^{-0.06}\cdot 0.34
    + e^{-0.06}\cdot 0.18
    + e^{-0.06}\cdot 0.22 \bigr)
  \;=\; 0.1144.
\end{equation}
This is exactly the Longstaff--Schwartz paper's headline result.

\subsection{Why in-the-money paths only?}

At $t=2$ in the worked example, OTM paths ($\Spx > \Kstrike$)
contribute zero to the put payoff because $(\Kstrike - \Spx)^+ = 0$
rules out immediate exercise.  Including them contaminates the fit:
the regression wastes capacity over a region where the exercise
decision is moot, hurting accuracy near the boundary where it matters.
The original paper proves restriction to ITM paths is asymptotically
consistent~\citep{longstaffschwartz2001}.

\subsection{Dimensionality: where LS genuinely wins}

For a single-stock American put, CRR trees with $N=5000$ steps
deliver four-decimal accuracy in $\sim 100\,\mathrm{ms}$ --- LS is
slower.  LS wins decisively once the state space has two or more
dimensions:
\begin{itemize}
  \item \textbf{Bermudan swaption} on a two-factor LMM model with
    40 exercise dates and a dozen forward rates;
  \item \textbf{Callable range accrual} whose coupon depends on the
    whole path of a reference rate;
  \item \textbf{Multi-asset American basket option} on ten stocks
    with a $10 \times 10$ correlation matrix;
  \item \textbf{American convertible bond} driven by stock, credit,
    and rates jointly.
\end{itemize}
The state $\Spx_t$ becomes a vector $\bm{X}_t$ and the basis becomes
a tensor product or low-order polynomial of its components; the rest
of the algorithm is unchanged.  A tree explodes as $2^d$ per step;
LS scales linearly in paths.  This is why LS displaced trees as the
industry default for exotic Americans.

\subsection{Code listing: Longstaff--Schwartz in $\sim$30 lines}

\begin{lstlisting}[caption={Longstaff--Schwartz American put pricer.}, label={lst:longstaff_schwartz}]
import numpy as np

def longstaff_schwartz_put(S0, K, r, sigma, T, M, N, seed=0):
    """Price an American put on GBM by Longstaff-Schwartz MC."""
    rng = np.random.default_rng(seed)
    dt = T / M
    df = np.exp(-r * dt)
    # Simulate GBM paths under Q (exact log-Euler).
    Z = rng.standard_normal((N, M))
    logS = np.log(S0) + np.cumsum(
        (r - 0.5*sigma**2)*dt + sigma*np.sqrt(dt)*Z, axis=1)
    S = np.hstack([np.full((N,1), S0), np.exp(logS)])  # shape (N, M+1)

    # Terminal payoff.
    CF = np.maximum(K - S[:, -1], 0.0)
    # Backward induction: for m = M-1 down to 1.
    for m in range(M-1, 0, -1):
        itm = S[:, m] < K
        if not itm.any():
            CF = CF * df
            continue
        X = S[itm, m]
        Y = CF[itm] * df   # realised discounted continuation
        A = np.column_stack([np.ones_like(X), X, X**2])
        beta, *_ = np.linalg.lstsq(A, Y, rcond=None)
        C_hat = A @ beta
        exer = K - X
        ex = exer > C_hat
        # Update cashflow for exercisers.
        CF_new = CF * df  # default: carry & discount
        idx = np.where(itm)[0][ex]
        CF_new[idx] = exer[ex]
        CF = CF_new
    return df * CF.mean()
\end{lstlisting}

Run on $\Spx_0 = 36$, $\Kstrike = 40$, $\rrf = 6\%$, $\volat = 20\%$,
$\Tmat = 1$ year, $M = 50$ exercise dates, $N = 100\,000$ paths
reproduces Longstaff and Schwartz's Table~1 price of $4.478$ to three
decimal places.

\begin{intuitionbox}
LS converts the optimal-control problem (when to exercise?) into a
sequence of regressions (what is the expected continuation value?).
The regressed $\hat C(\cdot)$ is biased --- only as good as the basis
--- but the resulting exercise policy is \emph{low-biased}: a
suboptimal rule gives a lower bound on the true American price,
useful for conservative pricing.  Upper bounds from duality methods
exist too, but LS low bounds are the industry workhorse.
\end{intuitionbox}

\subsection{Choosing the basis, and why the low-biased property matters}

The original paper uses weighted Laguerre polynomials $L_k(\Spx)$
of order two or three on ITM paths.  We used monomials above and
still reproduced the price because eight paths cannot distinguish
two bases of the same rank.  For larger $N$, Laguerre and Hermite
polynomials are better conditioned than raw monomials and should be
preferred once $\Spx$ spans more than an order of magnitude.  Use
three to five basis functions per state-variable
dimension~\citeGlasserman; beyond that, flexibility overfits.

The low-biased property matters.  LS's stopping rule is suboptimal
(the regression approximates), so the realised value $\hat
V_0^{\mathrm{LS}}$ along chosen stopping times is a lower bound on
the true American price.  Combined with a dual-martingale upper bound
(Rogers 2002, Andersen--Broadie 2004) one gets a tight confidence
interval.  In production the LS price is the point estimate; the
dual bound is spot-checked on regression-risk-sensitive products to
ensure the interval stays tight.

\begin{pitfallbox}[title={LS regressing all paths rather than ITM only}]
A common bug: regressing over all paths at each step instead of only
ITM ones.  OTM paths add noise where the exercise decision is moot
and degrade the fit near the exercise boundary where it matters.
Always filter to ITM paths before regressing.
\end{pitfallbox}


%% ================================================================
\section{Finite-difference PDE methods}
\label{sec:finite_difference}
%% ================================================================

Finite-difference methods solve the Black--Scholes PDE from
\cref{ch:black_scholes} \emph{directly on a grid}: discretise both
time and asset price, replace derivatives with differences, step
backward from the terminal payoff.  This is the natural tool when you
already have the PDE but no closed-form solution.  For low-dimensional
products (one or two state variables) with a PDE characterisation ---
vanilla European and American options, barriers, single-curve rates
products --- finite-difference methods win.  They are deterministic,
converge at $O(\Delta t + \Delta S^2)$ or better, and deliver
high-quality Greeks as grid by-products.

\subsection{The Black--Scholes PDE on a grid}

Recall from \cref{ch:black_scholes} that the value $V(t, \Spx)$ of a
European derivative on a non-dividend asset satisfies
\begin{equation}
\label{eq:bs_pde_ch34}
  \frac{\partial V}{\partial t}
  + \tfrac12 \volat^2 \Spx^2 \frac{\partial^2 V}{\partial \Spx^2}
  + \rrf \Spx \frac{\partial V}{\partial \Spx}
  - \rrf V \;=\; 0,
  \qquad V(\Tmat, \Spx) = g(\Spx).
\end{equation}
Discretise on $\Spx_i = i\Delta\Spx$, $i = 0,\ldots,I$, and $t_n =
T - n\Delta t$, $n = 0,\ldots,N$, with $V_i^n \approx V(t_n, \Spx_i)$.
Centred differences give
\begin{equation}
\label{eq:fd_coeffs}
  \alpha_i = \tfrac12 \Delta t \bigl(\volat^2 i^2 - \rrf i\bigr), \quad
  \beta_i  = -\Delta t \bigl(\volat^2 i^2 + \rrf\bigr), \quad
  \gamma_i = \tfrac12 \Delta t \bigl(\volat^2 i^2 + \rrf i\bigr),
\end{equation}
and the PDE becomes (after backward-in-time substitution
$\tau = T - t$)
\begin{equation}
  V_i^{n+1} - V_i^n \;\approx\;
    \alpha_i V_{i-1}^{?} + \beta_i V_i^{?} + \gamma_i V_{i+1}^{?},
\end{equation}
where the choice of $V_i^{?}$ (old time $n$, new time $n+1$, or
average) defines the scheme.

\subsection{Explicit scheme}

Choosing $V_i^{?} = V_i^n$ gives the \emph{explicit} update
\begin{equation}
  V_i^{n+1} \;=\; V_i^n + \alpha_i V_{i-1}^n
                + \beta_i V_i^n + \gamma_i V_{i+1}^n.
\end{equation}
It is trivial to implement (no linear solves) but suffers a CFL-type
stability restriction: one can show~\citeHull\ that stability requires
\begin{equation}
\label{eq:explicit_stab}
  \Delta t \;\lesssim\; \frac{(\Delta \Spx)^2}{\volat^2 \Spx_{\max}^2},
\end{equation}
so halving the space step forces a \emph{quartering} of the time
step.  Unusable for fine grids.

\textbf{Concrete illustration.}  For $\Spx_{\max} = 400$, $I = 1000$
(so $\Delta\Spx = 0.4$), $\volat=0.2$, the bound is $\Delta t \lesssim
2.5\times 10^{-5}$: pricing a $T=1$ option needs over $40\,000$ time
steps, and halving $\Delta\Spx$ quarters $\Delta t$.  Explicit schemes
are competitive only on coarse grids and rarely appear in production,
but they are pedagogically clean.

\subsection{Implicit scheme}

Choosing $V_i^{?} = V_i^{n+1}$ gives the \emph{implicit} update
\begin{equation}
  (1 - \beta_i) V_i^{n+1} - \alpha_i V_{i-1}^{n+1} - \gamma_i V_{i+1}^{n+1}
  \;=\; V_i^n,
\end{equation}
a tridiagonal system solved in $O(I)$ per step by Thomas's algorithm.
Unconditionally stable but first-order in time.

\subsection{Crank--Nicolson: the industry default}

\citet{cranknicolson1947} averaged explicit and implicit:
\begin{equation}
  (I - \tfrac12 L) V^{n+1}
  \;=\; (I + \tfrac12 L) V^n,
\end{equation}
where $L$ is the tridiagonal matrix with rows $(\alpha_i, \beta_i,
\gamma_i)$.  The resulting scheme is
\begin{itemize}
  \item unconditionally stable for any $\Delta t, \Delta \Spx$;
  \item second-order accurate in \emph{both} $\Delta t$ and $\Delta
    \Spx$ (locally $O(\Delta t^2 + \Delta \Spx^2)$);
  \item one tridiagonal solve per time step, i.e.\ as cheap as
    implicit;
  \item the canonical choice for single-factor European and American
    options in practice.
\end{itemize}

\begin{keyfactbox}[Crank--Nicolson summary]
\textbf{Crank--Nicolson} on the Black--Scholes PDE is
\emph{unconditionally stable} and \emph{second-order accurate in
space and time}.  It is the default scheme for any single-factor
equity or rates derivative that has a PDE characterisation.  For
American options, supplement each time step with an elementwise
$V_i^{n+1} \leftarrow \max(V_i^{n+1}, g(\Spx_i))$ projection to
enforce the exercise constraint.
\end{keyfactbox}

\subsection{Worked example: European call via Crank--Nicolson}

Same call ($\Spx_0=100$, $\Kstrike=100$, $\rrf=5\%$, $\volat=20\%$,
$\Tmat=1$), grid with $\Spx_{\max} = 400$, $I = 1000$, $N = 100$.
Terminal $V_i^0 = (\Spx_i - \Kstrike)^+$; boundary $V(0, t) = 0$ and
$V(\Spx_{\max}, t) \approx \Spx_{\max} - \Kstrike\, e^{-\rrf\tau}$.

Running Crank--Nicolson and interpolating at $\Spx_0 = 100$ gives
\begin{equation}
  V^{\mathrm{CN}}_0 = 10.4502, \qquad
  \text{vs.}\ \Call^{\mathrm{BS}} = 10.4506,
  \qquad |{\mathrm{err}}| = 3.8\times 10^{-4}.
\end{equation}
Four-decimal accuracy in a PC-second.

\subsection{Greeks for free from the grid}

A bonus of finite-difference: the full Greek surface falls out of the
grid without bumping.  Since the grid stores $V(t, \Spx_i)$ at every
node,
\begin{align}
  \Delta(t, \Spx_i) &\approx
    \frac{V(t, \Spx_{i+1}) - V(t, \Spx_{i-1})}{2\Delta\Spx}, \\
  \Gamma(t, \Spx_i) &\approx
    \frac{V(t, \Spx_{i+1}) - 2 V(t, \Spx_i) + V(t, \Spx_{i-1})}{(\Delta\Spx)^2}, \\
  \Theta(t, \Spx_i) &\approx
    \frac{V(t + \Delta t, \Spx_i) - V(t, \Spx_i)}{\Delta t},
\end{align}
all from a single grid solve.  MC needs one price per bumped parameter
per Greek --- a $2 N_\text{paths}$ penalty per Greek, mitigated only
by pathwise or likelihood-ratio estimators (\citeGlasserman, ch.~7).

\subsection{American options on the FD grid}

For an American option, add one line per step: after the linear solve,
project onto the exercise payoff,
\begin{equation}
  V_i^{n+1} \;\leftarrow\; \max\!\bigl( V_i^{n+1},\; g(\Spx_i) \bigr).
\end{equation}
This is the discrete analogue of the free-boundary condition
$V(t, \Spx) \geq g(\Spx)$ and delivers American prices at the same
$O(NI)$ cost.  The projection can be formulated as a linear-
complementarity problem (LCP) solved by projected SOR for higher
boundary accuracy; simple projection is fine for single-stock puts.

\begin{intuitionbox}
For a single-asset product with a PDE, finite difference annihilates
MC: an $O(IN)$ grid gives four decimals where $10^6$ MC paths give
two.  The dimension curse bites at 2--3 state variables; by 4--5
factors the tensor grid is dead and MC wins.
\end{intuitionbox}


%% ================================================================
\section{Trees: the CRR binomial model}
\label{sec:crr_trees}
%% ================================================================

Before FD libraries and Longstaff--Schwartz, American options were
priced on \emph{lattices}.  The \citet{coxrossrubinstein1979}
binomial tree (CRR) is the simplest, still taught everywhere and
occasionally used in production for simple contracts (small American
options, discretely-monitored single-asset barriers).  The idea:
approximate continuous-time lognormal GBM by a discrete up/down
lattice, which makes early-exercise decisions trivial to check at
each node (exactly what BS closed-form cannot do) and converges to
BS as the step count $N \to \infty$.

\subsection{Construction}

Discretise $[0, \Tmat]$ into $N$ steps of length $\Delta t = \Tmat/N$.
At each step the underlying goes up by $u$ or down by $d$:
\begin{equation}
  u = e^{\volat \sqrt{\Delta t}},
  \qquad d = 1/u = e^{-\volat \sqrt{\Delta t}},
  \qquad p = \frac{e^{\rrf \Delta t} - d}{u - d}.
\end{equation}
This matches the first two moments of $\log \Spx$ under $\Qm$ and
converges to GBM as $N \to \infty$.  The tree is \emph{recombining}
(up-down = down-up), giving $n+1$ nodes at time $n$, total $O(N^2)$.

\begin{figure}[ht]
\centering
\begin{tikzpicture}[
  nd/.style={circle, draw, minimum size=1.1cm, inner sep=1pt,
    font=\scriptsize, fill=blue!10, align=center},
  arr/.style={-Stealth, thin, gray!65},
]
% t=0
\node[nd] (s0)   at (0,0)    {$S_0$};
% t=1
\node[nd] (su)   at (2.5, 1.5) {$S_0 u$};
\node[nd] (sd)   at (2.5,-1.5) {$S_0 d$};
% t=2
\node[nd] (suu)  at (5, 3)   {$S_0 u^2$};
\node[nd] (sud)  at (5, 0)   {$S_0 ud$};
\node[nd] (sdd)  at (5,-3)   {$S_0 d^2$};
% t=3
\node[nd] (suuu) at (7.5, 4.5) {$S_0 u^3$};
\node[nd] (suud) at (7.5, 1.5) {$S_0 u^2d$};
\node[nd] (sudd) at (7.5,-1.5) {$S_0 ud^2$};
\node[nd] (sddd) at (7.5,-4.5) {$S_0 d^3$};
% t=0 to t=1
\draw[arr] (s0)  -- node[above left,  font=\tiny]{$p$}     (su);
\draw[arr] (s0)  -- node[below left,  font=\tiny]{$1\!-\!p$} (sd);
% t=1 to t=2
\draw[arr] (su)  -- (suu);  \draw[arr] (su)  -- (sud);
\draw[arr] (sd)  -- (sud);  \draw[arr] (sd)  -- (sdd);
% t=2 to t=3
\draw[arr] (suu) -- (suuu); \draw[arr] (suu) -- (suud);
\draw[arr] (sud) -- (suud); \draw[arr] (sud) -- (sudd);
\draw[arr] (sdd) -- (sudd); \draw[arr] (sdd) -- (sddd);
% Time labels
\foreach \t/\x in {0/0,1/2.5,2/5,3/7.5}
  \node[font=\footnotesize, below] at (\x,-5.4) {$t=\t$};
\node[font=\footnotesize, gray] at (1.25, 2.0) {$\times u$};
\node[font=\footnotesize, gray] at (1.25,-2.0) {$\times d$};
\end{tikzpicture}
\caption{Three-period CRR binomial tree. At each step the price moves up by
factor $u = e^{\sigma\sqrt{\Delta t}}$ or down by $d = 1/u$, with risk-neutral
probability $p = (e^{r\Delta t} - d)/(u - d)$.
The recombining property ($ud = du = S_0$) means the tree has $N+1$ terminal
nodes after $N$ steps, not $2^N$. Option prices are obtained by
\emph{backward induction}: compute payoffs at $t=3$, then take risk-neutral
expectations one step at a time back to $t=0$.}
\label{fig:binomial_tree}
\end{figure}

\subsection{Backward induction}

At expiry, the option value at each terminal node equals the payoff.
Moving back one step, the value at an internal node $(n, j)$ is
\begin{equation}
  V_{n,j} \;=\; e^{-\rrf \Delta t}
                \bigl[ p\, V_{n+1,j} + (1-p)\, V_{n+1,j+1} \bigr].
\end{equation}
For an American option, take $V_{n,j} = \max(g(\Spx_{n,j}), \text{the
above})$.  This is exactly \eqref{eq:dpeqn} executed on a tree.

\textbf{Convergence.}  For the same call parameters, CRR gives
\begin{center}
\begin{tabular}{r|rr}
\toprule
$N$ & CRR price & error vs BS \\
\midrule
10   & 10.2534 & $-0.1972$ \\
50   & 10.4107 & $-0.0399$ \\
200  & 10.4406 & $-0.0100$ \\
1000 & 10.4486 & $-0.0020$ \\
\bottomrule
\end{tabular}
\end{center}
Error scales as $1/N$, matching CRR's known $O(1/N)$
convergence~\citeHull.  For $10^{-4}$ accuracy, a $1000\times 100$
CN grid is about as fast as a CRR tree with $N=5000$.

\begin{intuitionbox}
CRR is pedagogically beautiful --- it visualises risk-neutral
valuation as a discrete conditional expectation --- and cheap for
simple single-factor Americans.  For serious work, Crank--Nicolson
(better accuracy, same cost) or Longstaff--Schwartz (path-dependence,
multi-factor) supersedes it.
\end{intuitionbox}

\subsection{Greeks from a CRR tree}

Like the FD grid, CRR delivers Greeks from stored node values:
\begin{equation}
  \Delta_0 \;\approx\; \frac{V_{1,0} - V_{1,1}}{\Spx_0(u-d)},
  \qquad
  \Gamma_0 \;\approx\; \frac{(V_{2,0}-V_{2,1})/(\Spx_0 u^2 - \Spx_0) - (V_{2,1}-V_{2,2})/(\Spx_0 - \Spx_0 d^2)}{\tfrac12 (\Spx_0 u^2 - \Spx_0 d^2)}.
\end{equation}
$\Theta$ needs two trees with different $\Tmat$ or a pairwise time-
differencing argument.  Vega and rho require a full rebuild with
bumped $\volat$ or $\rrf$, so tree Greeks are cheap for spot-dependent
quantities and expensive for others.

\subsection{Tree variants worth knowing}

\textbf{Black--Derman--Toy (BDT)} trees extend the binomial to short-
rate models with time-varying vol, still used for simple callable-bond
pricing.  \textbf{Trinomial trees} add a middle branch, eliminating
CRR's $O(\sqrt{\Delta t})$ odd/even oscillation near strike and
accelerating convergence to $O(1/N^2)$.  \textbf{Implied trees}
(Rubinstein 1994, Derman--Kani 1994) calibrate transition probabilities
from market option prices to construct a smile-consistent lattice;
they were the local-vol workhorse before Dupire's PDE superseded them.


%% ================================================================
\section{Greeks in Monte Carlo: pathwise and likelihood-ratio}
\label{sec:mc_greeks}
%% ================================================================

Unlike FD grids, Monte Carlo does not give Greeks for free.  The
naive \emph{bump and reprice} estimator
\begin{equation}
  \hat\Delta \;=\; \frac{\hat V_0(\Spx_0 + \epsilon) - \hat V_0(\Spx_0 - \epsilon)}{2\epsilon}
\end{equation}
has two problems: (i) fresh paths give independent noise with variance
$\propto 2/\sqrt N$, amplified by $1/\epsilon$; (ii) bump discretisation
error is $O(\epsilon^2)$.  The fix --- \emph{common random numbers}
across bumped prices --- kills independent noise and cuts Greek
variance by $10$--$100\times$.

Two more sophisticated estimators dominate production:

\textbf{Pathwise (adjoint) Greeks.}  If the payoff $g$ is almost-surely
differentiable in the parameter, one can swap differentiation and
expectation,
\begin{equation}
  \frac{\partial}{\partial \Spx_0}\, \E[g(\Spx_\Tmat(\Spx_0))]
  \;=\; \E\!\left[ \frac{\partial g}{\partial \Spx_\Tmat}
                   \cdot \frac{\partial \Spx_\Tmat}{\partial \Spx_0} \right].
\end{equation}
For GBM, $\partial \Spx_\Tmat / \partial \Spx_0 = \Spx_\Tmat / \Spx_0$,
and for a European call, $\partial g / \partial \Spx_\Tmat = \ind_{\Spx_\Tmat
> \Kstrike}$, so the pathwise delta is
\begin{equation}
  \hat \Delta^{\mathrm{pw}} \;=\; e^{-\rrf \Tmat}\, \frac{1}{N}
  \sum_{k=1}^N \ind_{\Spx_\Tmat^{(k)} > \Kstrike}\, \frac{\Spx_\Tmat^{(k)}}{\Spx_0}.
\end{equation}
No bump, no bias.  Pathwise fails when the payoff has a kink or jump
(digitals, barrier knock-outs): differentiating the indicator gives
a delta function an expectation cannot swallow.

\textbf{Likelihood-ratio (LR) Greeks.}  LR differentiates the density
instead of the payoff:
\begin{equation}
  \frac{\partial}{\partial \theta}\, \E[g(\Spx_\Tmat)]
  \;=\; \E\!\left[ g(\Spx_\Tmat)\, \frac{\partial \log p(\Spx_\Tmat; \theta)}{\partial \theta} \right].
\end{equation}
The score $\partial_\theta \log p$ is smooth, so LR works even for
discontinuous payoffs at the cost of higher variance for smooth ones.
The rule: \emph{pathwise for smooth, LR for discontinuous}; when both
apply, pathwise usually has smaller variance.  \citeGlasserman\
chapter~7 is canonical.

\begin{keyfactbox}[Monte Carlo Greek estimators]
\textbf{Bump-and-reprice with common random numbers} is the universal
default.  \textbf{Pathwise} is better when the payoff is
differentiable.  \textbf{Likelihood-ratio} is necessary when the
payoff has jumps (digitals, barriers).  A production Monte Carlo
engine typically implements all three, dispatched by payoff type.
\end{keyfactbox}


%% ================================================================
\section{Method selection: the decision table}
\label{sec:method_selection}
%% ================================================================

The flowchart below is what every derivatives-group Strat
internalises by the end of their first year.

\begin{center}
\small
\renewcommand{\arraystretch}{1.2}
\begin{tabular}{p{0.38\linewidth} | p{0.55\linewidth}}
\toprule
\textbf{Situation} & \textbf{Recommended method} \\
\midrule
Closed-form formula exists (vanilla European on BS, forward,
cash-and-carry bound, etc.).
& Use the closed form.  No numerical error, constant-time. \\
\midrule
European payoff, 1--2 underliers, PDE known.
& Crank--Nicolson finite difference.  Second-order accurate,
unconditionally stable, Greeks free from the grid. \\
\midrule
European payoff, $\geq 3$ underliers, or complicated path-dependence.
& Monte Carlo with antithetic + control variate.  Curse-of-dimension-
free. \\
\midrule
American or Bermudan exercise, 1--2 state variables.
& Crank--Nicolson with per-step projection onto the exercise payoff,
or CRR tree if you need to ship it this afternoon. \\
\midrule
American or Bermudan exercise with path-dependence, or $\geq 3$
factors.
& Longstaff--Schwartz Monte Carlo.  The only scalable option. \\
\midrule
Deep out-of-the-money / rare event (credit default, CVaR tail,
autocallable KO).
& MC with importance sampling.  Plain MC wastes $99\%$ of paths. \\
\midrule
Characteristic function known (Heston, VG, Kou, affine models).
& Fourier inversion (Carr--Madan FFT).  $O(N \log N)$ per strike
slice, closed-form accuracy. \\
\bottomrule
\end{tabular}
\end{center}

\begin{keyfactbox}[Numerical method decision tree]
\begin{enumerate}
  \item Closed form $\to$ use it.
  \item Low-dim European $\to$ Crank--Nicolson.
  \item High-dim European $\to$ MC with variance reduction.
  \item Low-dim American $\to$ CN with projection, or CRR.
  \item High-dim or path-dependent American $\to$
    Longstaff--Schwartz MC.
  \item Characteristic function known $\to$ Fourier.
\end{enumerate}
When in doubt, implement two methods and cross-check: agreement to
$10^{-4}$ is the standard test that neither pricer has a bug.
\end{keyfactbox}


%% ================================================================
\section{Back-references: where each method shows up}
\label{sec:method_backrefs}
%% ================================================================

Earlier Goldman-oriented chapters use these methods without naming
them:

\begin{itemize}
  \item \textbf{\Cref{ch:rates_derivatives} (rates).}
    Hull--White caps have a closed form under the HJM measure change;
    the Bermudan swaption alongside does not and uses Longstaff--
    Schwartz MC on the short-rate SDE.  Single-factor Hull--White is
    also cheap on a 1D Crank--Nicolson PDE in $r$, which gives
    crisper short-rate Greeks.
  \item \textbf{\Cref{ch:equity_derivatives} (equity).}
    Heston and rough-vol have known characteristic functions, so
    Carr--Madan FFT prices Europeans in milliseconds.  Path-dependent
    exotics (autocallables, lookbacks, cliquets) fall back on MC with
    antithetics, strata, and a BS control variate.
  \item \textbf{\Cref{ch:credit_derivatives} (credit).}
    Vanilla CDS prices are closed-form PVs of survival-weighted
    cashflows.  Gaussian-copula CDO tranches on 120 names are
    simulated by MC with importance-sampling tilt on the common
    factor~\citeGlasserman.
\end{itemize}

A derivatives book typically wraps all four pricers behind one
\verb|Pricer.price(product)| facade with a config flag.  The
engineering challenge is rarely the maths --- it is wiring the right
solver to the right product and making every solver return the same
Greek.

\subsection{Calibration: the outer loop}

Pricers sit inside a calibration loop that tunes the model to liquid-
vanilla market prices.  For local-vol:
\begin{equation}
  \min_{\volat_{\text{loc}}(t, \Spx)}
  \sum_i \left[ \Call^{\mathrm{model}}(\Kstrike_i, \Tmat_i;\, \volat_{\text{loc}})
                - \Call^{\mathrm{market}}(\Kstrike_i, \Tmat_i) \right]^2,
\end{equation}
with each $\Call^{\mathrm{model}}$ evaluation a PDE or MC run.  A
thousand quotes over a hundred LM iterations fires the inner pricer
$10^5$ times, so inner-pricer speed is the dominant library metric.
When accuracy ties, the faster method wins --- why Carr--Madan FFT
(Heston) and Crank--Nicolson (local-vol) dominate their niches.

\subsection{Numerical hazards in production}

A partial list of bugs every derivatives library has hit once:
\begin{itemize}
  \item \textbf{Off-by-one discount factor.}  One extra step between
    exercise and today costs $\sim 0.02\%$ per basis point, surfacing
    as systematic PV bias under \cref{ch:risk_production}'s P\&L-
    explain scrutiny.
  \item \textbf{Antithetic on two-sided kink payoffs.}  Variance
    silently increases, the price still converges --- just slower.
    Regress empirical SE against $1/\sqrt N$ as a check.
  \item \textbf{LS regression on OTM paths.}  Price biased low by
    $5$--$20$ bp depending on moneyness, never flagged by unit tests.
  \item \textbf{Explicit FD blowing up on a vol bump.}  User bumps
    vol from $20\%$ to $40\%$; \eqref{eq:explicit_stab} is violated;
    prices explode.  Fix: ban explicit schemes in production.
  \item \textbf{Basis conditioning.}  Raw monomials $\{1, \Spx,
    \Spx^2, \Spx^3\}$ give near-singular regression matrices when
    $\Spx$ has wide range; use orthogonal polynomials or rescale to
    $[-1, 1]$.
\end{itemize}


%% ================================================================
\section{Benchmarks and cross-checks}
\label{sec:benchmarks}
%% ================================================================

Every derivatives library ships a benchmark suite that cross-checks
each pricer against others at canonical parameter points:

\begin{enumerate}
  \item \textbf{Vanilla European call, flat BS.}  Every pricer matches
    \eqref{eq:rn_pricing} to $10^{-4}$ of notional, catching discount-
    factor and time-convention bugs.
  \item \textbf{Put--call parity.}  $\Call - \Put = \Spx - \Kstrike\,
    e^{-\rrf\Tmat}$ holds path-by-path on shared paths.
  \item \textbf{MC vs CN convergence curves.}  Error vs $1/\sqrt N$
    for MC, vs $\Delta t$ for CN; slope deviations flag bugs.
  \item \textbf{LS vs CRR for small-state American puts.}  A tree
    with $N=10^4$ gives $<10^{-3}$ accuracy as the LS reference.
  \item \textbf{Greek sanity.}  Grid $\Delta, \Gamma, \Theta$ agree
    with BS to $<1\%$ and with MC CRN bumps within CI.
\end{enumerate}

These run nightly in CI.  On a Strats desk the rule is: \emph{no
pricer ships without a reference pricer that agrees to tolerance at
three points}.  Choosing the right reference is often harder than
implementing the primary pricer.

\begin{keyfactbox}[Two-pricer cross-check]
Never trust a single pricer in production.  Every product family has
at least two (closed form + MC, CN + tree, or MC + LS); convergence
agreement at a benchmark point is the only reliable correctness test.
Unit tests without cross-pricer agreement are theatre.
\end{keyfactbox}


%% ================================================================
\section{Summary}
\label{sec:num_summary}
%% ================================================================

\begin{itemize}
  \item The toolkit has four members: closed form, Monte Carlo, finite
    difference, trees.  Pick whichever delivers required accuracy
    fastest given dimension and exercise style.
  \item \textbf{Monte Carlo}: $O(N^{-1/2})$, dimension-free, the only
    option for high-dim baskets or path-dependence.
  \item \textbf{Variance reduction} --- antithetics, control variates,
    importance sampling --- cuts MC variance by $5$--$100\times$.
  \item \textbf{Longstaff--Schwartz}: the canonical American pricer.
    Run paths forward; at each exercise date regress continuation on
    a polynomial basis.  $\sim$30 lines price the universe.
  \item \textbf{Crank--Nicolson}: default 1D--2D PDE scheme.
    Unconditionally stable, second-order accurate, linear-time per
    step.  A $100\times 1000$ grid matches BS to four decimals in
    a second.
  \item \textbf{CRR trees}: $O(1/N)$, easy to reason about, legitimate
    for very small American jobs.  Superseded by CN and LS otherwise.
  \item The flowchart in \S\ref{sec:method_selection} is what every
    derivatives Strat carries in their head.  When in doubt, cross-
    check two methods.
\end{itemize}


\begin{prosperitybox}
The Prosperity backtester is itself a Monte Carlo simulator: each day
replays historical order flow with pseudo-random counterparty
arrivals, and your PnL is one draw from a stochastic distribution.
Variance-reduction logic transfers: run many days (MC averaging),
compute paired deltas when you tweak a parameter (antithetic: same
shocks, new parameter), and use a simpler strategy as a control
variate.  The Frankfurt Hedgehogs credit ``backtest on as many days
as possible'' as their single most important rule --- variance
reduction applied to algo trading research.
\end{prosperitybox}


\begin{gsbox}
On a Strats desk the numerical toolkit is productised as a library
--- call it \verb|PricerLib| in SecDB dialect --- that all desks
consume.  A structuring Strat rarely writes MC from scratch: the
workflow is (i) check the product-to-pricer routing table, (ii) pick
the default (CN for low-dim exotics, LS for Americans, Carr--Madan
for Hestons), (iii) add a cross-check in the alternate pricer and
diff-test, (iv) hand off to the risk book.  The differentiator is
knowing which solver is numerically reliable for a product family
and owning the stability and convergence tests in CI.
\end{gsbox}


\begin{historybox}[title={Three papers that run derivatives desks}]
\begin{itemize}
  \item \citet{cranknicolson1947}: time-averaging scheme for the heat
    equation, Cambridge, decades before Black--Scholes.  Almost every
    parabolic PDE in finance is the heat equation in disguise.
  \item \citet{coxrossrubinstein1979}: binomial tree as a pedagogical
    device, which became the first tractable American pricer and the
    industry standard for 22 years.
  \item \citet{longstaffschwartz2001}: replaced the tree with a
    regression on simulated paths, breaking the curse of dimension
    for Americans.  The most-cited paper in computational finance ---
    read the eight-path Section~1 example until you can reproduce it
    without looking.
\end{itemize}
\end{historybox}
